\selectlanguage{ngerman}%
\frenchspacing%
\begin{ledgroupsized}[r]{120mm}%
\footnotesize%
\pstart%
\noindent\textbf{Überlieferung:}%
\pend%
\end{ledgroupsized}%
\begin{ledgroupsized}[r]{114mm}%
\footnotesize%
\pstart%
\parindent -6mm%
\makebox[6mm][l]{\textit{L}}%
Aufzeichnung:
LH~XXXVII,\;7 Bl.\;8.
Ein Blatt 4\textsuperscript{o},
Ränder beschnitten.
Zwei ganzseitig beschriebene Seiten.
\pend%
\end{ledgroupsized}%
\vspace{5mm}%
\begin{ledgroup}%
\footnotesize%
\pstart%
\noindent%
\textbf{Datierungsgründe:}
Die Datierung der Aufzeichnung N.\;\ref{42488} beruht auf %
inhaltlichen Betrachtungen.
Siehe die ausführliche Begründung auf S.\;\pageref{42488_dat}.%
\pend%
\end{ledgroup}%
\selectlanguage{latin}%
\frenchspacing%
\vspace{8mm}%
\normalsize%
\pstart%
\noindent%
\lbrack%
\href{\Scan{LH_37_07_008r}}%
{8\,r\textsuperscript{o}}%
\rbrack\
\edtext{%
Succi nutritii ascensus%
\indexS{succus nutritius in plantis}%
\indexS{ascensus succi in plantis}%
}{%
\lemma{Succi}%
\Bfootnote{%
\textit{(1)}~cum %
\textit{(2)}~nutritii ascensus%
~\textit{L}%
}}
in plantis%
\indexS{planta}
non est deducendus ex communi illa aeris pressione%
\indexS{pressio aeris communis}
quae Mercurium in Barometro%
\indexS{mercurius suspensus in barometro}%
\indexS{barometrum}
\edtext{%
ad 27 circiter%
}{%
\lemma{ad 27}%
\Bfootnote{%
\textit{(1)}~\angles{vel 28} %
\textit{(2)}~circiter%
~\textit{L}%
}}
pollices%
\indexS{pollex}
suspensum tenet,%
\indexS{suspensio mercurii in barometro}
aut aquam in antlia aspirante%
\indexS{aqua in antlia aspirante}%
\indexS{antlia aspirans}
ad 32 pedes%
\indexS{pes}
attollit,%
\indexS{suspensio aquae in antlia}
neque inde colligere licet
\edtext{%
arbores%
\indexS{arbor excrescens}
ultra 17 ulnarum%
\indexS{ulna}
altitudinem%
\indexS{altitudo arborum}
excrescere non debere,%
}{%
\lemma{arbores}%
\Bfootnote{%
\bars{non \textit{gestr.}} ultra 17 %
\textit{(1)}~ulnas excrescere %
\textit{(2)}~ulnarum altitudinem excrescere non debere,%
~\textit{L}%
}}
cum multo altiores pinus%
\indexS{pinus}
\edtext{%
abietesve%
\indexS{abies}
haberi%
}{%
\lemma{abietesve}%
\Bfootnote{%
\bars{apud nos \textit{gestr.}} haberi%
~\textit{L}%
}}
res sit
\edtext{%
indubitata;%
\indexS{res indubitata}
ut taceam%
}{%
\lemma{indubitata;}%
\Bfootnote{%
\textit{(1)}~sed re %
\textit{(2)}~ut taceam%
~\textit{L}%
}}
non omnem succum%
\indexS{succus nutritius in plantis}
inde a radice venire,%
\indexS{radix plantarum}
sed novum ex aere%
\indexS{aer pluvius}
aquave pluvia,%
\indexS{aqua pluvia}
imo etiam ex vicina
\edtext{%
planta sustentante interdum%
}{%
\lemma{planta}%
\Bfootnote{%
\textit{(1)}~interdum %
\textit{(2)}~sustentante interdum%
~\textit{L}%
}}
attrahi.%
\indexS{attractio succi in plantis}
\pend%
\pstart%
\edlabel{LH_37_07_008_jel-1}%
Succi igitur propagationem%
\indexS{propagatio succi in plantis}%
\indexS{succus nutritius in plantis}
rectius peti
putem ab alio naturae
\edtext{%
principio,%
\indexS{principium naturae}
nimirum a \textls{pressione subtiliore}%
\indexS{pressio aeris subtilioris}%
}{%
\lemma{principio,}%
\Bfootnote{%
\textit{(1)}~quod exp\textlangle on\textrangle\ %
\textit{(2)}~nimirum a \textls{pressione subtiliore}%
~\textit{L}%
}}
nondum satis explorata,
\edtext{%
a pondere%
\indexS{pondus aeris}
et vi elastica%
\indexS{vis elastica aeris}
\lbrack aeris\rbrack\
non dependente,%
\indexS{pressio aeris independens a pondere et vi elastica}%
}{%
\lemma{a pondere et vi elastica %
\bars{aeris \textit{erg. Hrsg. nach \refpassage{LH_37_07_008_008r_hdf-1}{LH_37_07_008_008r_hdf-2}}} %
non dependente,}%
\Bfootnote{%
\textit{erg.\;L}%
}}
quae multis tamen sese effectibus ostendit.%
\indexS{effectus pressionis aeris subtilioris}
Ex his unus
\edtext{%
est
cum%
}{%
\lemma{est}%
\Bfootnote{%
\textit{(1)}~quod %
\textit{(2)}~cum%
~\textit{L}%
}}
\edlabel{LH_37_07_008_008r_xyz_1}%
aquam%
\indexS{aqua assurgens sponte sua in tubis capillaribus}
vel alios multos liquores%
\indexS{liquor assurgens sponte in tubis capillaribus}
experimur
in tubos vitreos tenuissimos etiam capillares%
\indexS{tubus vitreus capillaris}
\edtext{%
(qualibus celeberrimus Leewenhoekius%
\indexN{\textls{Leeuwenhoek}, Antoni van 1632\textendash1723}
ad experimenta microscopica utitur)%
\indexS{experimentum microscopicum}%
}{%
\lemma{qualibus \lbrack...\rbrack\ utitur}%
\Cfootnote{%
Siehe \textsc{A.\;van Leeuwenhoek},
"Letter wherein some account is given of the manner of observing so great a number of little animals in divers sorts of water"
(\textit{PT} XII, Nr.\;134, 23.\;April [3.\;Mai] 1677, S.\;844).\cite{02520}
Der Kapillaritätseffekt wird dort nicht erwähnt.
Auf das Phänomen geht Leibniz aber ein, als er Ende 1677 in einem Brief für J.\;P. de La Roque\indexN{\textls{La Roque} (Larroque), Jean Paul de, gest. 1691} über Leeuwenhoeks\indexN{\textls{Leeuwenhoek}, Antoni van 1632\textendash1723} Mikroskopie berichtet (\textit{LSB} III,\;2 N.\;94, S.\;258.1--6, 261.7--10).\cite{02521}
Bei einem persönlichen Gespräch in Delft\indexO{Delft} mitte November 1676 (\textit{Chronik}, S.\;46)\cite{01236} hatte Leeuwenhoek\indexN{\textls{Leeuwenhoek}, Antoni van 1632\textendash1723} allerdings von seinen neuesten Forschungen erzählt;
vgl. Leibnizens Brief an H.\;Oldenburg\indexN{\textls{Oldenburg} (Grubendol), Heinrich 1618\textendash1677} vom 18./28. November 1676 (\textit{LSB} III,\;5 N.\;II, S.\;6.17--22).\cite{02522}
Leeuwenhoek\indexN{\textls{Leeuwenhoek}, Antoni van 1632\textendash1723} könnte hierbei auch über seine Erfahrungen mit dem Kapillaritätseffekt berichtet haben.%
}}
ad pollicum%
\indexS{pollex}
aliquot altitudinem%
\indexS{altitudo liquorum in tubis capillaribus}
sponte sua
\edtext{%
assurgere.%
\indexS{aqua assurgens sponte sua in tubis capillaribus}%
\indexS{liquor assurgens sponte in tubis capillaribus}
\edtext{%
Quod primus dicitur observasse Rohaultius,%
\indexN{\textls{Rohault}, Jacques 1618\textendash1672}
et extant complura de hac re experimenta%
\indexS{experimentum de liquoribus in tubis capillaribus}
in Monconisi\bracks{i}%
\indexN{\textls{Monconys} (Monconisius), Balthasar de 1611\textendash1665}
itinerario relata.%
}{%
\lemma{Quod \lbrack...\rbrack\ relata}%
\Cfootnote{%
Siehe \textsc{B.\;de Monconys},
\glqq Discours sur l’ascension de l’eau en un tuyau estroit\grqq\cite{02519}
(\textit{Journal de voyages}, Bd.\;III, Lyon 1666,
\glqq Lettres escrittes à Monconys\grqq, S.\;33--38).\cite{00118}%
}}%
\edlabel{LH_37_07_008_008r_xyz_2}
Hanc%
}{%
\lemma{assurgere.}%
\Bfootnote{%
\textit{(1)}~Hanc %
\textit{(2)}~\textlangle--\textrangle\ %
\textit{(3)}~Quod primus %
\lbrack...\rbrack\ experimenta in %
\bars{Monconisio \textit{ändert Hsrg.}} itinerario relata. Hanc%
~\textit{L}%
}}
pressionem%
\indexS{pressio aeris subtilioris}
etiam arbitror rotundare
\edtext{%
guttas%
\indexS{gutta rotundata}%
\indexS{rotundatio guttarum}
(usque adeo ut duae vicinae minores in unam novam coeant)%
\indexS{guttae coeuntes in unam}
et exprimere%
\indexS{gutta expressa}%
}{%
\lemma{guttas}%
\Bfootnote{%
\textit{(1)}~et exprimere etiam %
\textit{(2)}~(usque adeo %
\lbrack...\rbrack\ unam novam %
\textit{(a)}~cone %
\textit{(b)}~coeant) et exprimere%
~\textit{L}%
}}
rotundatas%
\indexS{gutta rotundata}
in eam partem
qua liquor facilius fertur;
veluti
\edtext{%
si atramenti guttam%
\indexS{gutta atramenti}%
\indexS{atramentum}
vel aliam superficiem tangamus%
\indexS{superficies tacta apice pennae}
apice pennae jam imbutae,%
\indexS{apex pennae}%
\indexS{penna imbuta atramento}%
}{%
\lemma{si}%
\Bfootnote{%
\textit{(1)}~guttam atramenti pennae jam imbutae tangamus %
\textit{(2)}~atramenti guttam %
\lbrack...\rbrack\ jam imbutae,%
~\textit{L}%
}}
videbimus atramentum%
\indexS{atramentum}
ex gutta%
\indexS{gutta atramenti}
in pennam%
\indexS{penna imbuta atramento}
subito ferri
sive sursum sive ad latus,
prout
\edlabel{LH_37_07_008_gravitasaetheris-1}%
\edtext{%
se offeret situs.%
\indexS{situs}
\edtext{%
Haec omnia compertum mihi est fieri
etiam in
\edtext{%
vacuo Gerikiano%
\indexS{vacuum Gerickianum}%
}{%
\lemma{vacuo Gerikiano}%
\Cfootnote{%
Künstlich erzeugtes Vakuum.
Anspielung auf den in
\textsc{O.\;v.\;Guericke}, \textit{Experimenta nova}, lib.\;II, cap.\;3 (Amsterdam 1672, S.\;56a)\cite{00055}
dargelegten Begriff der "Leere" sowie auf die
ebd., lib.\;III (S.\;71--124)\cite{00055}
beschriebenen Versuche hierüber.
}}
ex quo omnis pene aer est extractus%
\indexS{aer extractus}%
}{%
\lemma{Haec \lbrack...\rbrack\ extractus}%
\Cfootnote{%
Quelle nicht ermittelt.%
}}
\edlabel{LH_37_07_008_008r_hdf-1}%
atque adeo ponderi%
\indexS{pondus aeris}
vel vi elasticae aeris%
\indexS{vis elastica aeris}
tribui nequeunt.%
\edlabel{LH_37_07_008_008r_hdf-2}
\edtext{%
Quidam viri docti putant%
\bracks{,}%
\indexS{viri docti}
praeter aeris atmosphaeram,%
\indexS{atmosphaera aeris}%
\indexS{aer incumbens telluri}
incumbere telluri nostrae%
\indexS{tellus}
aliam atmosphaeram%
\indexS{atmosphaera alia}
multo altiorem%
\indexS{atmosphaera altior}
et subtiliorem%
\indexS{atmosphaera subtilior}
et ipsum aerem penetrantem%
\indexS{atmosphaera penetrans ipsum aerem}%
}{%
\lemma{Quidam \lbrack...\rbrack\ penetrantem}%
\Cfootnote{%
Siehe N.\;\ref{60205} in diesem Band, bes. %
\refpassage{LH_37_03_115-116_115_lpo-1}{LH_37_03_115-116_115_lpo-2}.%
}}
(ut aqua spongiam)%
\indexS{aqua penetrans spongiam}%
\indexS{spongia}
cujus pressione fiat%
\indexS{pressio aeris subtilioris}
ut
\edtext{%
liquor assurgat in Tubulum arctiorem utrinque apertum,%
\indexS{liquor assurgens in tubulum utrinque apertum}%
\indexS{tubulus apertus utrinque}%
}{%
\lemma{liquor \bracks{...} apertum}%
\Cfootnote{%
Siehe oben,
\refpassage{LH_37_07_008_008r_xyz_1}{LH_37_07_008_008r_xyz_2}.%
}}
et ut
\edtext{%
Mercurius aere purgatus%
\indexS{mercurius purgatus aere}
in Tubo superius
\lbrack%
\href{\Scan{LH_37_07_008v}}%
{8\,v\textsuperscript{o}}%
\rbrack\
clauso%
\indexS{tubus clausus superius}
ad 70 usque pollices%
\indexS{pollex}
suspensus teneatur%
\indexS{mercurius suspensus ad 70 pollices}%
\bracks{,}%
}{%
\lemma{Mercurius \bracks{...} teneatur}%
\Cfootnote{%
Hierüber berichtet
\textsc{J.\;Wallis}, \textit{Mechanica}, cap.\;XIV, prop.\;13, schol.
(London 1670--1671, Bd.\;II, S.\;740;\cite{00301}
\textit{WO}~I, S.\;1051);\cite{01008}
Leibniz hatte 1674--1675 die Stelle exzerpiert
(vgl. \textit{LSB} VIII,\;2 N.\;8, S.\;96.18--97.6).%
\cite{01343}
Siehe zudem
J.\;\textsc{Wallis}, "An extract of letters concerning the suspension of quicksilver well purged of air" %
(\textit{PT} VII, 1672, S.\;5160).\cite{00213}%
}}
et
\edtext{%
siphones%
\indexS{sipho}
in vacuo etiam Gerikiano%
\indexS{vacuum Gerickianum}
eundem Mercurium aere purgatum%
\indexS{mercurius purgatus aere}
attollant%
\bracks{,}%
}{%
\lemma{siphones \bracks{...} attollant}%
\Cfootnote{%
Wohl Anspielung auf
\textsc{C.\;Huygens}, "Extrait d'une lettre, touchant les phenomenes de l'eau purgée d'air"
(\textit{JS}, 25.\;Juli 1672, S.\;140;\cite{00062}
\textit{HO} VII, S.\;206).\cite{00113}
Im erwähnten (VI.) Experiment wurde aber Wasser verwendet, nicht Quecksilber.%
}}
et
\edtext{%
duo marmor\bracks{a} probe levigata%
\indexS{marmora levigata}
experiente Cl. Voldero
longe plus ponderis%
\indexS{pondus conans distrahere}
perpendiculariter distrahere conantis%
\indexS{pondus conans distrahere}
sustineant %
quam est incumbentis aeris cylindri.%
\indexS{cylinder aeris incumbentis}%
\indexS{aer incumbens}%
}{%
\lemma{duo \bracks{...} cylindri}%
\Cfootnote{%
Dieses im März 1676 von B.\;de Volder\indexN{\textls{Volder} (Volter, Wolder), Burchard (Burcher) de 1643\textendash1709} in Leyden\indexO{Leyden} durchgefürte Experiment kannte Leibniz wohl aus einem Bericht in der von P.\;Lohmeier betreuten Dissertation \textit{Exercitationum physicarum de paradoxi gravitatis et levitatis prima}, parad.\;I, §~IX (Rinteln 1678, S.\;32--34).\cite{01429}
Das Experiment erwähnt Lohmeier\indexN{\textls{Lohmeier} (Lohmeierus, Lomeierus), Philipp 1648\textendash1680} auch in seinem Brief an Leibniz vom 23.\;August (2.\;September) 1679
(\textit{LSB} III,\;2 N.\;336, S.\;817.16--20).\cite{01425}
Über einen ähnlichen Versuch de Volders\indexN{\textls{Volder} (Volter, Wolder), Burchard (Burcher) de 1643\textendash1709} berichtete ferner die anonyme Bemerkung "Experience faite à Leyden" im Heft des \textit{Journal des sçavans} vom 17.\;April 1679 (S.\;107\,f.).\cite{02524}
Hierauf spielt E.\,W.\,v.\;Tschirnhaus\indexN{\textls{Tschirnhaus} (H. v. Tch., H. v. Tsch.), Ehrenfried Walther v. 1651\textendash1708} in seinem Brief an Leibniz vom 11.\;September 1682 ab
(\textit{LSB} III,\;3 N.\;399, S.\;711.8\,f.).\cite{02525}
Siehe %
N.\;\ref{60205} in diesem Band,
\refpassage{LH_37_03_115-116_115r_deVolder-1}{LH_37_03_115-116_115v_deVolder-2}.%
}} %
Sed licet haec nova Atmosphaera subtilior%
\indexS{atmosphaera subtilior}
sufficere posset
reddendae rationi%
\indexS{ratio reddenda}
suspensi Mercurii aere purgati%
\indexS{mercurius purgatus aere}%
\indexS{mercurius suspensus ad 70 pollices}
et marmoreae Tabulae pondere gravatae,%
\indexS{tabulae marmoreae gravatae}
non tamen apparet
quid conferat ad liquorem impellendum%
\indexS{liquor impellendus}
in orificium inferius tubi%
\indexS{orificium tubi}
etiam supra aperti,%
\indexS{tubus apertus utrinque}
et multo minus
quomodo rotundet guttas%
\indexS{gutta rotundata}%
\indexS{rotundatio guttarum}
aut exprimat liquorem in locum%
\indexS{locus imbutus liquore}
cognato liquore jam imbutum.%
\indexS{liquor expressus in locum imbutum liquore}%
\edlabel{LH_37_07_008_jel-2}%
\edlabel{LH_37_07_008_gravitasaetheris-2}%
}{%
\lemma{se}%
\Bfootnote{%
\bars{penna \textit{streicht Hrsg.}} offeret %
\textit{(1)}~guttae. \lbrack/\rbrack\ %
\textit{(a)}~Hanc autem pressionem %
\textit{(b)}~Non est %
\bars{autem \textit{gestr.}} necessarium %
\textit{(c)}~Sed nec necessarium est ut haec pressio\indexS{pressio e gravitate aetheris} oriatur ex gravitate\indexS{gravitas aetheris} aetheris\indexS{aether} seu materiae aere subtilioris;\indexS{materia subtilior aere} cum videamus simplicem motum liquorum intestinum\indexS{motus intestinus liquorum} %
\bars{insensibilem \textit{erg.}}\indexS{motus insensibilis liquorum} rotundare guttas extranei liquoris, veluti olei in aqua natantis.\indexS{oleum natans in aqua} Itaque fieri posset ut Aethereae cujusdam materiae\indexS{materia aetherea} motus\indexS{motus materiae aethereae} haec efficeret, qui idem et vis Elasticae\indexS{vis elastica} causa\indexS{causa vis elasticae} %
\textit{(aa)}~esse \textlangle possit\textrangle\ %
\textit{(bb)}~est \textlangle--\textrangle\ %
\textit{(2)}~situs. %
\bars{Haec omnia %
\textit{(1)}~fiunt %
\textit{(2)}~compertum mihi %
\lbrack...\rbrack\ tribui nequeunt. \textit{erg.}} %
Quidam viri %
\lbrack...\rbrack\ aeris atmosphaeram, %
\textit{(a)}~occumbere %
\textit{(b)}~incumbere telluri %
\lbrack...\rbrack\ Tubulum arctiorem %
\textit{(aa)}~itemque ut Mercurius in %
\textit{(bb)}~utrinque apertum, et ut Mercurius %
\textit{(aaa)}~in Tubo %
\textit{(bbb)}~aere purgatus in Tubo superius %
\lbrack\href{\Scan{LH_37_07_008v}}{8\,v\textsuperscript{o}}\rbrack\ %
\textit{(aaaa)}~clausus %
\textit{(bbbb)}~clauso ad %
\lbrack...\rbrack\ et duo %
\bars{marmore \textit{ändert Hrsg.}} probe levigata %
\bracks{...} plus ponderis %
\bars{perpendiculariter distrahere conantis \textit{erg.}} %
sustineant, quam %
\lbrack...\rbrack\ cylindri. Sed %
\textit{(aaaaa)}~\textlangle-- --\textrangle\ %
\textit{(bbbbb)}~licet haec %
\lbrack...\rbrack\ conferat ad %
\textit{(aaaaa-a)}~tubum %
\textit{(bbbbb-b)}~liquorem impellendum %
\lbrack...\rbrack\ et multo minus %
\textit{(aaaaa-aa)}~\textlangle ponendo\textrangle\ %
\textit{(bbbbb-bb)}~quomodo rotundet %
\lbrack...\rbrack\ in locum %
\textit{(aaaaa-aaa)}~eodem %
\textit{(bbbbb-bbb)}~cognato liquore jam imbutum.%
~\textit{L}%
}}
\pend%
\pstart%
\edlabel{LH_37_07_008_008v_kex-1}%
Itaque malim
\edtext{%
esse Motum quendam%
}{%
\lemma{esse}%
\Bfootnote{%
\textit{(1)}~modum quendam %
\textit{(2)}~Motum quendam%
~\textit{L}%
}}
magnetico analogum%
\indexS{motus magneticus}%
\indexS{motus analogus magnetico}
in omni corporum sensibilium particula%
\indexS{partculae corporis sensibilis}%
\indexS{corpus sensibile}
tanquam ipsa magnes esset,%
\indexS{magnes}
et quamlibet quidem a qualibet
\edtext{%
nonnihil,
sed tamen magis cognatam a cognata trahi.%
\indexS{particulae cognatae}
Sane videmus magnetem armatum%
\indexS{magnes armatus}
non nisi accedente contactu%
\indexS{contactus}
vim majorem exercere%
\indexS{vis magnetica}%
\lbrack,\rbrack\
in distans vero non agere fortius
quam non armatum;
et in universum
aliqua a contactu adhaerescentia%
\indexS{adhaerescentia a contactu}
est%
}{%
\lemma{nonnihil,}%
\Bfootnote{%
\textit{(1)}~sed magis tamen cognatam trahi a cognata %
\textit{(a)}~(unde %
\textit{(aa)}~\textlangle--\textrangle\ %
\textit{(bb)}~ex contactu \textlangle--\textrangle\ %
\textit{(b)}~(unde ex quovis contactu adhaerentia quaedam %
\textit{(2)}~sed tamen %
\lbrack...\rbrack\ videmus magnetem %
\textit{(a)}~ope contactu %
\textit{(b)}~armatum non %
\lbrack...\rbrack\ majorem exercere %
\bars{in distans %
\bracks{...} non armatum \textit{erg.}}; %
\textit{(aa)}~ut adeo intelligi possit %
\textit{(bb)}~et in %
\lbrack...\rbrack\ adhaerescentia est%
~\textit{L}%
}}
liquorum ad liquores,%
\indexS{adhaerescentia liquorum ad liquores}
liquorum ad solida%
\indexS{adhaerescentia liquorum ad solida}
et solidorum inter se.%
\indexS{adhaerescentia solidorum inter se}
\edtext{%
Quanquam non dubitem
contactus loco%
\indexS{locus contactus}%
\indexS{contactus}%
}{%
\lemma{Quanquam}%
\Bfootnote{%
\textit{(1)}~\angles{--} %
\textit{(2)}~in %
\textit{(3)}~non dubitem %
contactus loco%
~\textit{L}%
}}
esse et viciniam insensibilis intervalli.%
\indexS{vicinia intervalli insensibilis}%
\indexS{intervallum insensibile}
Hujus contactus major vis est,%
\indexS{vis contactus}%
\indexS{contactus}
\edtext{%
cum nullus%
}{%
\lemma{cum}%
\Bfootnote{%
\textit{(1)}~non %
\textit{(2)}~nullus%
~\textit{L}%
}}
liquor heterogeneus%
\indexS{liquor heterogeneus}
intercedit,%
\indexS{liquor intercedens}
is enim magnam contactuum partem intercipit;%
\indexS{contactus}
inde Mercurius aere purgatus%
\indexS{mercurius purgatus aere}
partes habet magis haerentes inter se,%
\indexS{partes mercurii haerentes}
et magis ipse adhaeret%
\indexS{mercurius adhaerens vitro}
\edtext{%
vitro.%
\indexS{vitrum}
Et praeterea pressio ambientis%
\indexS{pressio ambientis}
resistit separationi%
\indexS{resistentia separationi}%
\indexS{separatio}
cum ambiens non admittitur in intervallum.%
\indexS{ambiens non admissum in intervallum}%
}{%
\lemma{vitro.}%
\Bfootnote{%
\textit{(1)}~Ut taci %
\textit{(2)}~Et praeterea pressio %
\textit{(a)}~vicinae si locus non commode a\textlangle--\textrangle\ %
\textit{(b)}~ambientis resistit separationi %
\textit{(aa)}~in quam %
\textit{(bb)}~cum ambiens %
\lbrack...\rbrack\ in intervallum.%
~\textit{L}%
}}
Sed ab hac sola pressione%
\indexS{pressio ambientis}
guttarum formatio%
\indexS{formatio guttarum}
et
\edtext{%
liquoris in tubos arctos sublevatio%
\indexS{sublevatio liquoris in tubos arctos}%
\indexS{tubus arctus}
peti nequit;%
}{%
\lemma{liquoris}%
\Bfootnote{%
\textit{(1)}~expressio peti nequit %
\textit{(2)}~in tubos %
\lbrack...\rbrack\ peti nequit;%
~\textit{L}%
}}
ita enim non tantum sustinetur aliquid,
sed etiam movetur.
\ledpb%
Itaque attractione opus est
quae Magneticae%
\indexS{attractio magnetica}
vel etiam Electricae%
\indexS{attractio electrica}
sit analoga,%
\indexS{attractio analoga}
quae sola fortasse differunt subtilitate%
\indexS{subtilitas}
et inter se%
\bracks{,}
et
subtiliore adhuc nostra.%
\indexS{attractio subtilior}%
\edlabel{LH_37_07_008_008v_kex-2}%
\pend%
\pstart%
Instituenda autem sunt experimenta.%
\indexS{experimentum instituendum}
Ex.\;gr. an Mercurius purgatus ab aere,%
\indexS{mercurius purgatus aere}
suspensus in tubo 70 pollicum%
\indexS{mercurius suspensus ad 70 pollices}
\edtext{%
cadat in vacuo%
\indexS{vacuum}
si vel%
}{%
\lemma{cadat}%
\Bfootnote{%
\textit{(1)}~, si ve %
\textit{(2)}~in %
\textit{(a)}~\lbrack p\rbrack\ %
\textit{(b)}~in vacuo si vel%
~\textit{L}%
}}
minimum foramen%
\indexS{foramen minimum}
in summitate aperiatur;
ita enim res non attribuenda
\pgrk{t\~h| sunafe'ia|}%
\indexS{\pgrk{sun'afeia}}
et motui ejus autori,%
\indexS{motus autor \pgrk{t\~hs sunafe'ias}}
sed pressioni ambientis.%
\indexS{pressio ambientis}
\edtext{%
I\bracks{t}em%
}{%
\lemma{Idem}%
\Bfootnote{%
\textit{L\;ändert Hrsg.}%
}}
erit
an experimentum marmorum politorum%
\indexS{experimentum marmorum politorum}%
\indexS{marmora polita}
effici etiam possit,
si tantum circumferentiae marmorum sint politae,
ad excludendum scilicet externum affluxum,
accedente et materia pingui.%
\indexS{materia pinguis}
Nam si res pendet a pressione ambientis,%
\indexS{pressio ambientis}
non a contactu%
\indexS{contactus}%
\bracks{,}
non ideo minus succedet experimentum.%
\indexS{experimentum marmorum politorum}
Itaque distinguendam arbitror
\edtext{%
adhaesionem%
\indexS{adhaesio}
a pressione ortam%
\indexS{pressio orta ab attractione}
ab ea quae est attractione
tum sese contingentium,%
\indexS{attractio contingentium sese}
tum etiam distantium.%
\indexS{attractio distantium}
Haec causa
quae guttas rotundat%
\indexS{causa rotundans guttas}%
\indexS{gutta rotundata}%
\indexS{rotundatio guttarum}
et liquorem in arctissimos tubos attollit.%
\indexS{causa attollens liquorem in tubos arctos}%
\indexS{liquor sublevatus in tubos arctos}%
\indexS{sublevatio liquoris in tubos arctos}%
}{%
\lemma{adhaesionem}%
\Bfootnote{%
\textit{(1)}~et attractionem %
\textit{(2)}~a pressione %
\lbrack...\rbrack\ est attractione %
\textit{(a)}~quae \textlangle eadem\textrangle\ et guttas rotundat %
\textit{(b)}~quae %
\textit{(c)}~tum sese %
\lbrack...\rbrack\ tubos attollit.%
~\textit{L}%
}}
\pend%
\pstart%
Cum fluidis tenuioribus corpora crassiora pervadantur,%
\indexS{fluidum tenue pervadens corpus crassius}%
\indexS{corpus crassum pervasum fuido tenuiori}
fluida autem illa sint in motu,%
\indexS{fluidum in motu}
hinc ubique
inter obstacula%
\indexS{obstaculum impediens motum rectilineum}
\edtext{%
motum rectilineum impeditum%
\indexS{motus rectilineus impeditus}
convertunt in circularem,%
\indexS{motus rectilineus conversus in circularem}%
\indexS{fluidum convertens motum rectilineum in circularem}%
\indexS{motus circularis e rectilineo}%
}{%
\lemma{motum}%
\Bfootnote{%
\textit{(1)}~\textlangle --\textrangle\ impeditum convertunt in rec %
\textit{(2)}~rectilineum impeditum convertunt in circularem,%
~\textit{L}%
}}
hinc motus vorticosus partium crassiorum fluidi interlabentis,%
\indexS{motus vorticosus partium crassarum fluidi}%
\indexS{partes crassae fluidi}%
\indexS{fluidum interlabens}
tum tenuiores%
\indexS{partes tenues fluidi}
facilius perlabuntur%
\indexS{partes fluidi perlabentes}%
\bracks{;}
porro ex motu vorticoso%
\indexS{motus vorticosus partium crassarum fluidi}
oritur vis expansiva,%
\indexS{vis expansiva orta e motu vorticoso}
oritur et perpetua
\angles{emissio}%
\indexS{emissio perpetua}
et
\edtext{%
eradiatio%
\indexS{eradiatio}
dum bullae rumpuntur.%
\indexS{bulla rupta}
Hinc autem%
}{%
\lemma{eradiatio}%
\Bfootnote{%
\textit{(1)}~; ex qua \textlangle--\textrangle\  %
\textit{(2)}~dum bullae rumpuntur. Hinc autem%
~\textit{L}%
}}
gravitas%
\indexS{gravitas orta a motu vorticoso}
et lux.%
\indexS{lux orta a motu vorticoso}
Etsi gravitatis vis%
\indexS{vis gravitatis}
penetrantior
quam lucis.%
\indexS{vis lucis}
\pend%
\vspace{3.0em}%
\pstart%
\noindent%
\centering%
*     *     *
\pend%
\newpage%
\selectlanguage{german}%
\frenchspacing%
\footnotesize%
\pstart%
\noindent%
\textls{Ausführliche Datierungsgründe:}%
\pagelabel{42488_dat}
Die titellose und undatierte Aufzeichnung N.\;\ref{42488} ist auf einem Träger überliefert, der weder Spuren eines Wasserzeichens noch sonstige äußere Merkmale aufweist, die eine eindeutige chronologische Einordnung ermöglichen würden. %
Die vorgeschlagene Datierung -- ab Sommer 1690 (unter Vorbehalt) -- lässt sich folgendermaßen begründen:
\pend%
\pstart%
(1)
Die Aufzeichnung N.\;\ref{42488} wurde mit großer Wahrscheinlichket nach der wohl zwischen Herbst 1680 und April 1683 abgefassten Aufzeichnung \textit{De nova pressione aeris subtilioris} (N.\;\ref{60205} im vorliegenden Band) angefertigt.
Die These eines subtileren Fluidums, dessen Druck die Phänomene -- zunächst Adhäsionsphänomene im Vakuum -- erklären sollte, welche nicht auf den Luftdruck zurückgeführt werden konnten, ist im ersten Teil von N.\;\ref{42488} (\refpassage{LH_37_07_008_jel-1}{LH_37_07_008_jel-2}) in kürzerer Form dargestellt als im umfangreicheren Text N.\;\ref{60205}, der sich ausführlicher mit den Einzelheiten der damaligen Forschung auseinandersetzt.
Daher die Annhme, dass die Aufzeichnung N.\;\ref{42488} nach \textit{De nova pressione aeris subtilioris} entstanden ist.
\pend%
\pstart%
(2)
Einige Aspekte von N.\;\ref{42488} könnten als Zeichen dafür gedeutet werden, dass der Text nicht viel später als N.\;\ref{60205} verfasst wurde, d.\,h. noch in den 1680er Jahren.
Das trifft zunächst auf das für N.\;\ref{42488} geradezu ausschlaggebende Thema der Kapillarität zu, das Leibniz wahrscheinlich schon seit der Pariser\indexO{Paris} Zeit bekannt war (siehe N.\;\ref{56216} im vorliegenden Band).
Auch die Quellen, die im ersten Teil von N.\;\ref{42488} (\refpassage{LH_37_07_008_jel-1}{LH_37_07_008_jel-2}) erwähnt werden, stammen aus den 1670er (oder gar 1660er) Jahren.
Ebenso verdächtig ist, dass die thematisch treffende Abhandlung \textit{De gravitate aetheris}\cite{01403} von Jacob Bernoulli,\indexN{\textls{Bernoulli}, Jacob 1655\textendash1705} die 1683 veröffentlicht wurde, in N.\;\ref{42488} nicht -- wie es zu erwarten wäre -- genannt wird.
(Der Begriff der \textit{gravitas aetheris} wird in N.\;\ref{42488} ausdrücklich in einer gestrichenen Variante verwendet; vgl. den textkritischen Apparat zu \refpassage{LH_37_07_008_gravitasaetheris-1}{LH_37_07_008_gravitasaetheris-2}.)
Demzufolge könnte man eine Datierung der Aufzeichnung N.\;\ref{42488} bis ungefähr 1686 für plausibel halten.
Dagegen spricht aber folgender Grund:
\pend%
\pstart%
(3)
Im zweiten Teil von N.\;\ref{42488} (\refpassage{LH_37_07_008_008v_kex-1}{LH_37_07_008_008v_kex-2}) lehnt Leibniz den Druck eines dünneren Fluidums als Ursache der Kapillarität ab und schlägt stattdessen eine auf der "Anziehung" (\textit{attractio}) beruhenden Erklärung vor, die sich nach Vorbildern wie Magnetismus und Elektrizität richten sollte.
Dieser Aspekt fehlte in der Aufzeichnung \textit{De nova pressione aeris subtilioris}, in der lediglich Phänomene, die entweder vom gemeinen Lufdruck oder vom Druck eines dünneren Fluidums herrühren sollten, in Betracht gezogen werden.
Die Unterscheidung zwischen \textit{pressio ambientis} und \textit{attractio} als Ursachen der Adhäsion in N.\;\ref{42488} ist ebenfalls neu und kann als Zeichen angesehen werden, dass die Aufzeichnung nicht vor Leibnizens Auseinandersetzung mit Newtons\indexN{\textls{Newton} (Neutonus), Isaac 1643\textendash1727} \textit{Principia}\cite{00535} %
und somit nicht älter sein dürfte als etwa der Text \textit{LSB} VI,\;4 N.\;380,\cite{02616} editorisch \textit{De firmitate, vi elastica, explosione, attractione} benannt.
In diesem auf die Monate zwischen Frühjahr und Sommer 1690 (unter Vorbehalt) datierbaren, sicherlich nach der \textit{Dynamica}\cite{01354} entstandenen Entwurf wird die "Anziehung" -- wohl zum ersten Mal bei Leibniz -- als Ursache der Kohäsion angeführt (vgl. \textit{LSB} VI,\;4 N.\;380, S.\;2081.4--7).\cite{02616}
Die Verwandschaft der Aufzeichnung N.\;\ref{42488} mit \textit{De firmitate, vi elastica, explosione, attractione} in der explanatorischen Rolle, die der \textit{attractio} jeweils zugewiesen wird, ist Grund für die Annahme, dass auch N.\;\ref{42488} nicht vor dem Frühjahr oder Sommer 1690 entstand.
\pend%
\pstart%
(4)
Eine spätere Entstehung der Aufzeichnung N.\;\ref{42488} ist nicht auszuschließen, da auch in späteren Texten die \textit{attractio} eine entscheidende Rolle bei der Erklärung der Kohäsion spielt, wie dies etwa in einer kurzen Aufzeichnung, die zwischen April 1704 und Anfang 1705\;(?) angefertigt wurde, der Fall ist (vgl. \textit{LSB} VIII,\;3 N.\;33, S.\;362.28--30).\cite{02617}
Einen Terminus ante quem der Datierung von N.\;\ref{42488} lässt sich jedoch nach heutigem Wissensstand nicht angeben.
\pend%
\pstart%
Die vorliegende Aufzeichnung N.\;\ref{42488} kann aus den angeführten Gründen (3) und (4) \textit{ab dem Frühjahr 1690} datiert werden.
Eine frühere Entstehung ist zwar weniger wahrscheinlich, aber aus den Gründen (1) und (2) nicht auszuschließen.
\pend%
\selectlanguage{latin}%
\frenchspacing%
\normalsize%
